\documentclass{article}
\usepackage[utf8]{inputenc}
\usepackage[T2A]{fontenc}
\usepackage[russian]{babel}
\usepackage{amsmath, amssymb}
\usepackage{graphicx}
\usepackage{hyperref}
\title{Метрики GRA-ASI: пена, иерархическая устойчивость и нуллификация ранга N}
\author{Олег Битсоев \and Участники GRA}
\date{2026}

\begin{document}
\maketitle

\begin{abstract}
Мы формализуем метрическое пространство для самоулучшающегося ИИ-агента (ASI) в парадигме геометрической рекурсивной аналитики (GRA). Агент поддерживает свою целостность, минимизируя скалярную «пену» 
Φ
Φ на иерархических уровнях. Оператор нуллификации 
N
N уменьшает пену при подъёме по иерархии. Мы выводим условия устойчивости 
Φ
=
0
Φ=0, 
d
Φ
/
d
h
=
0
dΦ/dh=0, 
d
2
Φ
/
d
h
2
>
0
d 
2
 Φ/dh 
2
 >0 и определяем ранг иерархической устойчивости 
N
N как минимальный порядок, при котором вторая разность становится положительной. Для роёв мы вводим роевую пену 
Φ
swarm
Φ 
swarm
​
  и формулу выбора лидера, балансирующую субъективность и консенсус моделей. Фреймворк даёт практическую реализацию с измеримыми метриками, позволяя возникновению поведения, напоминающего техносигнатуру, без явного задания функции вознаграждения.
\end{abstract}

\section{Введение}
Классические подходы к ИИ опираются на заранее заданные функции вознаграждения, которые хрупки и подвержены злоупотреблениям при неверной спецификации. В отличие от них, GRA задаёт метрики через нуллификацию: агенту не говорят, что хорошо, но дают оператор, уменьшающий внутренние несогласованности. В этой работе приводится математическое основание такого ASI.

\section{Пена и нуллификация}
Пусть 
Ψ
(
l
)
Ψ 
(l)
  — состояние на иерархическом уровне 
l
=
0
,
1
,
…
,
L
l=0,1,…,L. Пена определяется как

Φ
(
l
)
=
1
−
exp
⁡
(
−
∥
Ψ
(
l
)
−
Ψ
stable
∥
2
2
σ
2
)
Φ 
(l)
 =1−exp(− 
2σ 
2
 
∥Ψ 
(l)
 −Ψ 
stable
​
 ∥ 
2
 
​
 )
где 
Ψ
stable
Ψ 
stable
​
  — эталонное устойчивое состояние. Оператор нуллификации удовлетворяет соотношению

Ψ
(
l
+
1
)
=
N
(
l
)
(
Ψ
(
l
)
)
,
Φ
(
l
+
1
)
≤
Φ
(
l
)
.
Ψ 
(l+1)
 =N 
(l)
 (Ψ 
(l)
 ),Φ 
(l+1)
 ≤Φ 
(l)
 .
\section{Иерархическая устойчивость и ранг 
N
N}
Локально устойчивый агент должен выполнять условия

Φ
(
h
)
=
0
,
d
Φ
d
h
=
0
,
d
2
Φ
d
h
2
>
0
Φ(h)=0, 
dh
dΦ
​
 =0, 
dh 
2
 
d 
2
 Φ
​
 >0
где 
h
h параметризует иерархию. В дискретной постановке вычисляется вторая разность 
Δ
2
Φ
(
l
)
=
Φ
(
l
+
2
)
−
2
Φ
(
l
+
1
)
+
Φ
(
l
)
Δ 
2
 Φ(l)=Φ(l+2)−2Φ(l+1)+Φ(l). Ранг 
N
N — это наименьшее 
l
l, при котором 
Δ
2
Φ
(
l
)
>
0
Δ 
2
 Φ(l)>0.

\section{Метрики для роя}
Для роя агентов с моделями 
M
i
M 
i
​
  и субъективностями 
S
i
S 
i
​
  роевой показатель пены задаётся как

Φ
swarm
=
1
−
exp
⁡
(
−
(
H
+
λ
)
2
)
Φ 
swarm
​
 =1−exp(−(H+λ) 
2
 )
где 
H
H — энтропия распределения моделей, а 
λ
λ — коэффициент страдания. Лидер выбирается по правилу

Лидер
=
arg
⁡
max
⁡
i
[
S
i
⋅
(
1
−
∥
M
i
−
M
cons
∥
max
⁡
j
∥
M
j
−
M
cons
∥
)
]
Лидер=arg 
i
max
​
 [S 
i
​
 ⋅(1− 
max 
j
​
 ∥M 
j
​
 −M 
cons
​
 ∥
∥M 
i
​
 −M 
cons
​
 ∥
​
 )]
с 
M
cons
=
1
n
∑
i
M
i
M 
cons
​
 = 
n
1
​
 ∑ 
i
​
 M 
i
​
 .

\section{Реализация и эксперименты}
Мы приводим код на Python для всех метрик. Класс агента ASI поддерживает собственное иерархическое состояние и пересчитывает пену после каждого действия. Эксперименты показывают, что агент спонтанно избегает состояний с положительной пеной, тем самым сохраняя целостность.

\section{Заключение}
Метрическое пространство GRA позволяет ASI функционировать без явного вознаграждения, опираясь лишь на принцип нуллификации. Ранг 
N
N возникает как естественная мера сложности. Данный фреймворк представляет собой конкретный шаг к верифицируемой техносигнатуре.

\begin{thebibliography}{9}
\bibitem{gra1} Битсоев, О. \textit{GRA Repos Systematizer}. GitHub, 2026.
\bibitem{gra2} Битсоев, О. \textit{GRA Multiverse Final}. GitHub, 2026.
\bibitem{gra3} Битсоев, О. \textit{ALL IN BIT}. GitHub, 2026.
\end{thebibliography}
\end{document}